% ML-only paper — Phase 1 + Phase 2 (notebook: Projet7_brainTumorDetection.ipynb)
Brain tumor classification from magnetic resonance imaging (MRI) supports triage of glioma, meningioma, and pituitary lesions.
Most studies apply default parameters to handcrafted texture descriptors or couple exhaustive grid search to a single classifier, introducing optimisation bias and high compute cost.
We present a two-phase framework on a balanced dataset of 1{,}500 axial MRI slices (500 per class): Phase~1 selects GLCM, LBP, DWT, and HOG hyperparameters using three classifier-free separability metrics (Fisher discriminant ratio, mutual information, Davies--Bouldin index) over 213 configurations; Phase~2 benchmarks eight classical models (SVM, KNN, random forest, XGBoost, LightGBM, logistic regression, MLP, ExtraTrees) with stratified 5-fold \texttt{GridSearchCV} on eight optimised feature compositions.
On the held-out test split, the best configuration pairs KNN with the Full(opt) vector (accuracy $=0.953$, macro-F1 $=0.946$); SVM on Full(opt) reaches F1 $=0.912$, providing an interpretable clinical baseline.
Random-forest group importance confirms that HOG dominates raw importance while GLCM and statistical descriptors carry higher information per dimension, agreeing with Phase~1 separability rankings.
Full statistical validation (Cohen's $\kappa$, bootstrap confidence intervals, McNemar, Wilcoxon, Friedman--Nemenyi) is reported for the top models.
These results show that separability-guided handcrafted features with rigorous classical benchmarking achieve strong MRI classification without deep learning or GPU training.
